\section{Categorical Uniqueness of TtS}
\label{sec:g1-g6-categorical-uniqueness}

The exhaustion-by-deduction proof of \S\ref{sec:g1-g6-exhaustion-by-deduction} characterizes the trace-to-scheme operation as a morphism in $\mathbf{TtS}$ and derives sufficiency of $G_1$--$G_6$ from the morphism-data structure. The proof depends on the \emph{categorical-formalization assumption}: that $\mathbf{TtS}$ as defined captures the operation correctly.

This section closes that assumption by characterizing $\mathbf{TtS}$ uniquely up to categorical equivalence among categories that capture the trace-to-scheme operation. The characterization proceeds via three properties applied to morphism data: \emph{categorical necessity} (any operation-capturing category must include the six morphism-data pieces), \emph{categorical sufficiency} (a category with exactly these pieces captures the operation), and \emph{categorical minimality} (no proper subset of the pieces suffices). Together, these characterize $\mathbf{TtS}$ uniquely.

\subsection{Operation-capturing categories}
\label{subsec:operation-capturing}

\begin{definition}[Operation-Capturing Category]
\label{def:operation-capturing}
A category $\mathbf{C}$ \emph{captures the trace-to-scheme operation} if:
\begin{enumerate}[(i)]
\item Its objects include the $A_1$-admissible interfaces $\mathcal{T}, \mathcal{S}_1, \mathcal{S}_2, \dots$ as defined in \S\ref{subsec:TtS-category}.
\item Its morphisms include admissible export operations $\mathcal{T} \to \mathcal{S}_i$ for each scheme target.
\item Its morphism specification preserves enough structural data to determine which export is admissible (i.e., to determine whether $A_1$'s cost-payment requirements are met at each structural surface).
\end{enumerate}
\end{definition}

The third condition is the operative one: $\mathbf{C}$ must specify enough morphism data that admissibility can be checked. We now show this requires exactly six pieces of data.

\subsection{Categorical necessity of the six morphism-data pieces}
\label{subsec:cat-necessity}

\begin{lemma}[Categorical Necessity]
\label{lem:cat-necessity}
Any category $\mathbf{C}$ that captures the trace-to-scheme operation must include all six morphism-data pieces of $\mathbf{TtS}$ (target object, map graph, parameter set, probabilistic structure, role separation, comparison residual).
\end{lemma}

\begin{proof}
Suppose $\mathbf{C}$ omits one of the six pieces. Then there is a structural surface of the operation whose admissibility is unspecified in $\mathbf{C}$'s morphism data. Specifically:
\begin{itemize}
\item Omitting target object $\Rightarrow$ codomain undetermined; admissibility cannot be checked at the inter-category jump (cf. Lemma~\ref{lem:G1}).
\item Omitting map graph $\Rightarrow$ transport undetermined; admissibility cannot be checked at the value-transfer surface (cf. Lemma~\ref{lem:G2}).
\item Omitting parameter set $\Rightarrow$ external content undetermined; admissibility cannot be checked at the input-declaration surface (cf. Lemma~\ref{lem:G3}).
\item Omitting probabilistic structure $\Rightarrow$ uncertainty propagation undetermined; admissibility cannot be checked at the spread surface (cf. Lemma~\ref{lem:G4}).
\item Omitting role separation $\Rightarrow$ input/output mixing undetermined; admissibility cannot be checked at the smuggling surface (cf. Lemma~\ref{lem:G5}).
\item Omitting comparison residual $\Rightarrow$ post-comparison structure undetermined; admissibility cannot be checked at the residual surface (cf. Lemma~\ref{lem:G6}).
\end{itemize}
In each case, $\mathbf{C}$ fails Definition~\ref{def:operation-capturing}(iii): admissibility cannot be checked. Therefore $\mathbf{C}$ does not capture the operation. The contrapositive: $\mathbf{C}$ must include all six pieces. $\square$
\end{proof}

\subsection{Categorical sufficiency of the six morphism-data pieces}
\label{subsec:cat-sufficiency}

\begin{lemma}[Categorical Sufficiency]
\label{lem:cat-sufficiency}
A category $\mathbf{C}$ whose morphism specification consists of exactly the six morphism-data pieces captures the trace-to-scheme operation.
\end{lemma}

\begin{proof}
By Lemma~\ref{lem:morphism-data} of v35.4, the six pieces exhaust the morphism's data: any structural feature of an export operation corresponds to one of these pieces. By Theorem~\ref{thm:A1-forced} of v35.4, $A_1$ imposes cost-payment at each piece in bijection with $G_1$--$G_6$.

If $\mathbf{C}$'s morphism specification consists of these six pieces, then for any morphism in $\mathbf{C}$ the six structural surfaces are specified. Admissibility (whether $G_1$--$G_6$ close) can be checked surface-by-surface against the morphism's data. Therefore $\mathbf{C}$ captures the operation per Definition~\ref{def:operation-capturing}(iii). $\square$
\end{proof}

\subsection{Categorical minimality of the six morphism-data pieces}
\label{subsec:cat-minimality}

\begin{lemma}[Categorical Minimality]
\label{lem:cat-minimality}
No proper subset of the six morphism-data pieces captures the trace-to-scheme operation.
\end{lemma}

\begin{proof}
Suppose $\mathbf{C}$ has only five of the six pieces. By Lemma~\ref{lem:cat-necessity}, $\mathbf{C}$ does not capture the operation. The same argument extends to any proper subset.

Each missing piece has a witness from the framework's empirical training set:
\begin{itemize}
\item Missing target object: any unnamed-codomain comparison (general witness).
\item Missing map graph: routes that compare $T$ to $M$ without specifying $\tau$ (general witness).
\item Missing parameter set: routes using external constants without declaration (general witness).
\item Missing probabilistic structure: zero-$\sigma$ reports on transported value with positive input $\sigma$ (general witness).
\item Missing role separation: v18.0 W-route numeric projection (target consumed); v27 charged-lepton three-mode reconstruction (parameters consumed); v32 light-quark uniform scaling (PDG values consumed). All three quarantined under $G_5$ (v35.2 + v35.3).
\item Missing comparison residual: $H_0$ Route-V (residual unassignable to any channel; falsifying state under $G_6$).
\end{itemize}

Each piece is independently necessary; the six-piece set is minimal. $\square$
\end{proof}

\subsection{The categorical equivalence theorem}
\label{subsec:cat-equivalence}

\begin{theorem}[Categorical Uniqueness of $\mathbf{TtS}$]
\label{thm:cat-uniqueness}
Any category $\mathbf{C}$ satisfying categorical necessity (Lemma~\ref{lem:cat-necessity}), sufficiency (Lemma~\ref{lem:cat-sufficiency}), and minimality (Lemma~\ref{lem:cat-minimality}) is categorically equivalent to $\mathbf{TtS}$.
\end{theorem}

\begin{proof}
Let $\mathbf{C}$ be such a category. Define a functor $F: \mathbf{TtS} \to \mathbf{C}$ by:
\begin{itemize}
\item On objects: $F$ sends each $A_1$-admissible interface in $\mathbf{TtS}$ to the corresponding interface in $\mathbf{C}$ (which exists by Definition~\ref{def:operation-capturing}(i)).
\item On morphisms: $F$ sends each $\mathbf{TtS}$-morphism, specified by its six data pieces, to the $\mathbf{C}$-morphism with the matching six data pieces (which exists by categorical sufficiency).
\end{itemize}

\paragraph{$F$ is faithful.} Distinct $\mathbf{TtS}$-morphisms have distinct data tuples (target / map / parameters / probability / roles / residual). $F$ preserves the data tuples; therefore distinct $\mathbf{TtS}$-morphisms have distinct $\mathbf{C}$-images.

\paragraph{$F$ is full.} Every $\mathbf{C}$-morphism between $A_1$-admissible interfaces has six morphism-data pieces (by categorical necessity in $\mathbf{C}$ + minimality, $\mathbf{C}$'s morphism specification is exactly the six-piece structure). The $\mathbf{TtS}$-morphism with the same six pieces is the preimage.

\paragraph{$F$ is essentially surjective.} Every $\mathbf{C}$-object that is an $A_1$-admissible interface has a preimage in $\mathbf{TtS}$ by Definition~\ref{def:operation-capturing}(i). $\mathbf{C}$-objects that are not $A_1$-admissible interfaces are outside the operation-capturing scope and do not affect equivalence on the operation-capturing subcategory.

By the standard categorical-equivalence theorem (a fully faithful, essentially surjective functor establishes an equivalence of categories), $\mathbf{C}$ and $\mathbf{TtS}$ are categorically equivalent on their operation-capturing subcategories. $\square$
\end{proof}

\subsection{Joint statement of the deductive gate framework}
\label{subsec:deductive-complete}

Combining the structural-theorem stack of v35.2--v35.5:

\begin{theorem}[Trace-to-Scheme Export, Deductively Complete]
\label{thm:export-deductive}
A trace-to-scheme route admits physical export if and only if the six gates $G_1$--$G_6$ all close in their appropriate (numerical or structural) interpretation. The gate set is uniquely characterized by:
\begin{enumerate}[(a)]
\item \emph{Necessity}: each gate is $A_1$-derivable (Lemmas~\ref{lem:G1}--\ref{lem:G6} of v35.2).
\item \emph{Sufficiency} (two routes): closing all six gates suffices, by structural-exhaustion of operation surfaces (v35.3) or by morphism-data correspondence in $\mathbf{TtS}$ (v35.4 Theorem~\ref{thm:deductive-sufficiency}).
\item \emph{Minimality}: no proper subset of $G_1$--$G_6$ suffices (Corollary~\ref{cor:minimality} of v35.3).
\item \emph{Categorical uniqueness}: the trace-to-scheme category $\mathbf{TtS}$ is uniquely characterized up to equivalence among operation-capturing categories (Theorem~\ref{thm:cat-uniqueness} above).
\end{enumerate}
\end{theorem}

The framework's audit-first character is now deductively bound at the deepest level: $A_1$ requires exactly $G_1$--$G_6$, the gates cover exactly the morphism data of the categorically-unique $\mathbf{TtS}$, and the morphism data exhausts the operation's structural surfaces. There is no room for a hidden gate, no missing structural feature, no alternative formalization that would produce a different gate set.

\subsection{What's open after v35.5}
\label{subsec:post-v35-5}

The gate framework is deductively complete modulo two genuinely open programs:

\paragraph{Uniqueness of route-status determination.} The lemma that each gate-availability pattern leads to exactly one named status (export-candidate / external-MSR / obstruction-named / trace-with-residual / etc.) is empirically nine-for-nine across the day's articulated routes but not deductively proved. The proof would derive uniqueness from the gate-pattern combinatorics: enumerate possible gate-availability patterns; show each maps to exactly one status; show no pattern maps to two.

\paragraph{Conditionality on $A_1$.} The gate framework is uniquely characterized given $A_1$ as the foundational axiom. A different foundational axiom would produce a different category, different gates, and different export discipline. This is a feature, not a bug --- $A_1$ is the framework's load-bearing commitment, and the gate framework is the structural form of that commitment applied to the trace-to-scheme operation. But the gate framework cannot be reduced further than $A_1$; it inherits $A_1$'s status as a foundational assumption.

\paragraph{Practical content.} With v35.5 landed, the trace-to-scheme export theorem is publishable as a stand-alone structural-theorem paper independent of any per-route physics claim. The seven-layer artifact stack (v9.6 schema architecture; v35 export theorem; v35.1 structural extension; v35.2 necessity; v35.3 sufficiency-by-enumeration + minimality; v35.4 sufficiency-by-deduction; v35.5 categorical uniqueness) covers architecture, theorem, extension, necessity, sufficiency (twice), minimality, and uniqueness --- the full deductive structure of a structural-theorem program. A focused paper writing the v35--v35.5 chain end-to-end would be a substantial methodological contribution beyond APF's specific physics claims.
